%% Vélo dans une côte, de profil. Le repère est INCLINÉ : il suit la pente, comme celui qu'on
%% attache à la route et non à la feuille. y monte le long de la pente, x lui est perpendiculaire,
%% z sort de la feuille. La roue avant ① tourne autour de son axe (sortant), le vélo entier ②
%% avance le long de la pente — ni l'un ni l'autre ne se lit sans regarder l'inclinaison.
\documentclass[tikz,border=4pt]{standalone}
\usepackage{liaisons}
%% _objets-common.tex — outils communs aux figures « objet du quotidien avec son repère ».
%%
%% Ces figures servent des exercices où l'axe ne doit PAS être donné dans l'énoncé : l'élève lit
%% l'orientation du repère sur le dessin. D'où deux exigences :
%%   - le repère porte ses TROIS axes, dans une orientation qui change d'une figure à l'autre
%%     (si le vertical était toujours z, « vertical → z » redeviendrait un réflexe) ;
%%   - tout est en \large : à 390 px de large sur un téléphone, \small est illisible.
%%
%% Le \repere de liaisons.sty ne convient pas ici : deux axes seulement, et en \small.
%%
%% RÈGLE DE COULEUR, et elle est stricte : la couleur dit **mobile ou fixe** par rapport au
%% référentiel de l'exercice, jamais « la pièce dont on parle ». Le bâti — sol, route, mur, meuble,
%% tapis, plateau, support — est en solideB (bleu) ou en gris ; TOUT ce qui peut bouger par rapport
%% à lui est en solideA (rouge), même quand plusieurs pièces bougent. Ce sont les pastilles ① ② qui
%% désignent la pièce interrogée, c'est leur seule raison d'être. Colorer « la pièce ① » en rouge et
%% « la pièce ② » en bleu induirait en erreur une fois sur deux l'élève qui a appris que le rouge
%% bouge — le cas s'est produit sur le vélo, dont le cadre était bleu alors qu'il roule.

%% \repereXYZ[longueur]{point}{angle/label, angle/label, angle/label}
%% Trois flèches depuis un même point, chacune à l'angle voulu. L'axe « qui sort de la feuille »
%% se dessine comme les autres, en diagonale, à la manière des perspectives du cours.
\newcommand\repereXYZ[3][0.95]{%
  \foreach \ang/\lab in {#3} {%
    \draw[-{Stealth[length=6pt]}, line width=0.8pt] #2 -- ++(\ang:#1)
      node[pos=1, anchor={\ang+180}, inner sep=3pt, font=\large] {$\vec{\lab}$};%
  }}

%% \numero{point du repère}{point visé sur le dessin}{n} — pastille noire reliée à la pièce.
\newcommand\numero[3]{%
  \draw[line width=0.5pt, black!55] #1 -- #2;
  \node[circle, fill=black, text=white, inner sep=2pt, font=\large\bfseries] at #1 {#3};}

%% Épaisseurs : le dessin doit tenir à la taille d'un téléphone, donc peu de traits et des traits gras.
\tikzset{
  objet/piece/.style={line width=1.6pt, line cap=round, line join=round, draw=solideA},
  objet/bati/.style={line width=1.6pt, line cap=round, line join=round, draw=solideB},
  objet/fin/.style={line width=0.7pt, line cap=round, draw=black!55},
}

\begin{document}
\begin{tikzpicture}[x=1cm,y=1cm, rotate=18]

% --- la route : le bâti, en gris, pour que le bleu reste au vélo et le rouge à la roue
\draw[line width=1.6pt, line cap=round, black!55] (-0.7,-0.05) -- (6.2,-0.05);
\foreach \h in {-0.45,0.05,...,5.9} { \draw[objet/fin] (\h,-0.05) -- ++(-0.3,-0.32); }

% --- les deux roues
\def\rr{0.78}
\coordinate (av) at (4.5,\rr);
\coordinate (ar) at (1.3,\rr);
\draw[objet/piece, fill=white] (av) circle (\rr);
\draw[objet/piece, fill=white] (ar) circle (\rr);
\foreach \a in {0,60,120} {
  \draw[objet/fin] ($(av)+(\a:\rr)$) -- ($(av)+(\a+180:\rr)$);
  \draw[objet/fin] ($(ar)+(\a:\rr)$) -- ($(ar)+(\a+180:\rr)$);
}
\fill[solideA] (av) circle (0.11);
\fill[solideA] (ar) circle (0.11);

% --- le cadre : mobile par rapport à la route, donc rouge comme les roues.
% La couleur dit mobile ou fixe, pas « la pièce dont on parle » : ce sont les pastilles qui désignent.
\coordinate (ped) at (2.75,0.72);
\coordinate (sel) at (2.15,2.05);
\coordinate (gui) at (3.95,2.15);
\draw[objet/piece] (ar) -- (sel) -- (ped) -- cycle;
\draw[objet/piece] (sel) -- (gui) -- (ped);
\draw[objet/piece] (gui) -- (av);
\draw[objet/piece] ($(sel)+(-0.32,0.16)$) -- ($(sel)+(0.3,0.16)$);
\draw[objet/piece] ($(gui)+(-0.28,0.2)$) -- ($(gui)+(0.28,0.2)$);

% --- repère incliné : il suit la pente (le dessin entier est déjà tourné de 18°)
\repereXYZ{(0.15,3.15)}{90/y, 0/x, 210/z}

\numero{(5.35,2.4)}{($(av)+(35:\rr)$)}{1}
\numero{(3.0,3.5)}{(sel)}{2}   % à l'écart du repère, qu'elle recouvrait

\end{tikzpicture}
\end{document}
